\chapter{Experiments and Discussion}
\label{ch:experiments_discussion}

% ==============================================================================
% Core Research Question:
% How do the proposed Predictors and MPC formulations perform in closed-loop against
% the ground-truth Plant, what are their operational trade-offs, and where do they break down?
% ==============================================================================

% Contents:
% - Experimental Protocol and Evaluation Criteria:
%   * In-silico patient scenarios: Spontaneous seizure generation in EZ and recruitment into PZ.
%   * Quantitative metrics: Seizure State, Seizure Burden, total injected charge, solve times.
%   * Uncontrolled baseline and reference conditions.
% - Study 1: Offline Predictor Benchmarking:
%   * Selection of STFT geometry: Trading temporal responsiveness against frequency resolution.
%   * Training loss comparison: Multi-step Rollout stability of spectral training loss vs. raw MSE loss.
%   * Model comparison: Linear DMDc vs. Nonlinear MLP over the Control Horizon.
% - Study 2: Numerical Solver and Transcription Benchmarking:
%   * Solve time, convergence reliability, and constraint satisfaction across horizon lengths.
%   * Single shooting vs. multiple shooting performance.
% - Study 3: Closed-Loop Seizure Suppression:
%   * Proof of closed-loop efficacy: Attenuating seizure amplitude and preventing PZ recruitment via scalp EEG feedback.
%   * Waveform MPC vs. Spectral MPC: Empirical demonstration of why waveform tracking fails while spectral MPC succeeds.
%   * Continuous vs. Pulse-Based Stimulation: Trade-offs between Seizure Burden reduction, stimulation duration,
%     pulse counts, and solver complexity.
% - Critical Discussion & Failure Analysis:
%   * Analysis of negative results and controller breakdown modes (stochastic noise bursts, leadfield errors, latency limits).
%   * Translational implications of findings.
